\chapter*{Preface}

In 1948 a thirty-two-year-old engineer at Bell Telephone Laboratories published a paper that began with one sentence of
definition and ended with a new science. Claude Shannon showed that information can be measured, in bits, and that the
measure answers two practical questions exactly: how short a message can be made, and how fast it can be sent through a
noisy wire without error. The second answer was a shock. Noise, he proved, need not cause errors at all; it only lowers
the rate at which one can talk. He did not say how to do it, and finding codes that come close to his limits took the
next fifty years.

This book tells that story and two others that grew beside it. The first four chapters are about information:
entropy and its relatives, compression, the capacity of noisy channels, and Kolmogorov's measure of the information
in a single string, which explains why random data cannot be compressed and gives a strange new proof that there are
infinitely many primes. The next three are about codes: linear codes, whose most perfect example, Hamming's, hides the
exceptional lattice $E_8$ drawn on the cover; algebraic codes built from polynomials over finite fields, which let a
QR code survive a torn corner; and the codes on sparse graphs that finally reached Shannon's limits and now carry
mobile phone traffic. The last five are about secrets: perfect secrecy and why it is impractical, pseudorandomness
built from one-way functions, public-key cryptography and the number theory that attacks it, lattices and the
cryptography meant to survive quantum computers, zero-knowledge proofs and secure computation, and finally quantum
computing itself, with Shor's factoring algorithm and Grover's search.

The prerequisites are Tomes~1 to~4 of the main series and Tomes~C1 and~C2, with a few results from Tome~C3; they are
recalled where they are used, with a precise reference. Every theorem is proved, except for a few landmarks whose
proofs would need a book of their own; those are clearly labelled as quoted, with sources. Each chapter opens with a
question to think about, and ``Pause and think'' boxes ask you to predict before reading on. Exercises are graded from
$\star$ to $\star\star\star$; each has a hint and a solution that begins by explaining how one might find the idea.

Each chapter has a companion Python notebook that makes the mathematics run: Huffman and arithmetic coders, random
codes approaching capacity, Reed--Solomon and list decoders, a peeling decoder for sparse codes, attacks on toy
cryptosystems, the LLL algorithm breaking a knapsack, a small learning-with-errors scheme, secret sharing among three
parties, and a quantum computer simulated on a laptop, factoring $15$ and $21$. The theorems say what is possible and
what is not; the notebooks let you watch codes and ciphers come close to the limits, and fail beyond them.
